% Section 7.4, from lectures/L07/README.md: what you should be able to do, the questions to test
% yourself with, and the next lecture.

\section{Sammanfattning}
\label{c7:sec:summary}

\needspace{6\baselineskip}
\subsection*{Det här ska du kunna nu}
Efter det här kapitlet ska du kunna:
\begin{itemize}
  \item Beskriva ATmega328P:s delar, register, statusregister, programräknare och minnen, och
    koppla dem till mikrodatorns blockschema.
  \item Säga vilka register \code{ldi} kan ladda, och förklara varför utifrån instruktionens
    kodning.
  \item Förklara skillnaden mellan programminnet och dataminnet.
  \item Läsa en rad assembler och peka ut etikett, instruktion, operander och kommentar.
  \item Skriva samma tal decimalt, binärt och hexadecimalt i ett program.
  \item Skapa ett projekt i Microchip Studio, bygga det, rätta felen, och stega programmet i
    simulatorn.
  \item Förutsäga registrens värden efter varje instruktion i ett kort program, och kontrollera
    förutsägelsen i simulatorn.
\end{itemize}

\needspace{6\baselineskip}
\subsection*{Frågor att testa dig själv med}
\begin{enumerate}
  \item Varför kan \code{ldi} ladda \code{r16} men inte \code{r15}? Hur löser du det om du ändå
    behöver en konstant i \code{r5}?
  \item Den gula pilen i simulatorn står på \code{inc r18}. Har \code{inc r18} körts?
  \item Vad händer om du tar bort raden \code{rjmp end} sist i programmet?
  \item \code{ldi r16, 200}, \code{ldi r16, 0b11001000} och \code{ldi r16, 0xC8}: vad skiljer
    maskinkoden åt?
  \item Vad är skillnaden mellan en instruktion och ett direktiv? Blir \code{.def counter = r16}
    någon maskinkod?
  \item Ett program assembleras utan fel men ger fel resultat. Vad säger det om vad assemblern
    kontrollerar?
\end{enumerate}

\needspace{6\baselineskip}
\subsection*{Nästa kapitel}
Kapitel~\ref{c8:ch} handlar om aritmetik, flaggor och logik i registren:
\begin{itemize}
  \item Addition och subtraktion i registren, och vad som händer när resultatet inte får plats.
  \item Tvåkomplement: hur samma byte kan vara 255 eller $-1$.
  \item Flaggorna i SREG: vad varje instruktion lämnar efter sig.
  \item Logiska operationer, skift och rotation.
  \item Utporten: de första lysdioderna, i simulatorns I/O-fönster.
\end{itemize}
